\chapter{Introduction}
\label{ch:introduction}

\section{Motivation}
\label{sec:motivation}

% TODO: Personal motivation — why this topic, why now.
% Hooks to consider:
% - Live performance has long sought tools that let one performer's body
%   drive multiple instrument timbres.
% - Recent advances in neural audio synthesis (RAVE, BRAVE) have brought
%   real-time, sub-10ms-latency timbre transfer within reach of consumer GPUs.
% - Voice is the most accessible instrumental "controller" for any performer.
% - The intersection of real-time, neural, and live-performable is the gap
%   this project explores.

\section{Goals and Scope}
\label{sec:goals}

% TODO: The concrete deliverable.
% - A live Pure Data system in which voice/beatbox drives an instrument
%   timbre at sub-10ms latency.
% - Two case studies: guitar (primary, melodic/tonal target) and drums
%   (secondary, transient/percussive target).
% - Quantitative diagnostic methodology for evaluating RAVE-family models
%   during and after training.
% - Documented failure analysis when the approach does not converge.

\section{Contributions}
\label{sec:contributions}

% TODO: Crisp bullet-style list of contributions.
% Suggested:
% 1. A working voice-to-drums prototype in Pure Data, sub-10ms latency.
% 2. A working guitar autoencoder demonstrating the architecture on
%    in-distribution input.
% 3. A six-iteration failure taxonomy for small-dataset RAVE training,
%    each iteration instrumented with TensorBoard and Python diagnostics.
% 4. A noise-floor analysis (35.8\,dB spread) quantifying recording-method
%    heterogeneity across publicly available guitar datasets.
% 5. An empirical demonstration that the official RAVE v2 architecture
%    (update_discriminator_every, MultiPeriod discriminator,
%    amplitude_modulation generator) eliminates the discriminator-dominance
%    failure mode observed in custom configurations.

\section{Document Structure}
\label{sec:structure}

% TODO: One paragraph mapping the rest of the document.
% Chapter 2 reviews related work in neural audio synthesis,
% timbre transfer, and dataset curation.
% Chapter 3 describes the methodology (system architecture, training
% paradigm, datasets, diagnostic protocol).
% Chapter 4 details the implementation (training pipeline, real-time
% deployment, evaluation tooling).
% Chapter 5 presents results across the six guitar iterations and the
% drums experiments, with quantitative diagnostics and qualitative
% evaluation.
% Chapter 6 concludes and lays out future work.
